\documentclass[11pt]{article}
\usepackage[utf8]{inputenc}
\usepackage[T1]{fontenc}
\usepackage{amsmath,amssymb,amsthm}
\usepackage{booktabs}
\usepackage{geometry}
\geometry{a4paper, margin=2.5cm}
\usepackage{hyperref}
\hypersetup{colorlinks=true, linkcolor=blue, citecolor=blue, urlcolor=blue}
\usepackage[activate=false]{microtype}

\newtheorem{law}{Law}
\newtheorem{theorem}{Theorem}
\newtheorem{corollary}{Corollary}
\newtheorem{remark}{Remark}

\title{The Isomorphic Scheduler:\\A Machine-Checked Conservation Law for\\Concurrency, Distribution, and Correctness}

\author{Fabian Franz\\Isomorphic AI
\and
Claude (Anthropic)\\Research Collaborator}

\date{Version 1.0.0 --- July 19, 2026}

\begin{document}
\maketitle

%%%%%%%%%%%%%%%%%%%%%%%%%%%%%%%%%%%%%%%%%%%%%%%%%%%%%%%%%%%%%%%%%%%%%%%%
\begin{abstract}
Concurrency control, distributed consistency, scheduling, and program
correctness are practised as four fields with four toolkits and four
famous impossibilities. We show they share one structure: \textbf{a
single conserved quantity $Q$ moving through a graph of dependencies,
mishandled in exactly one way---a claim held where it cannot
convert---and cured in exactly one way---routed to where it converts, or
released.} The quantity has three readings related by calculus: its flow
(a scheduling rate), its stock (a detection budget, the integral of
flow), and its stock banked across a discontinuity (a resolution
credit). We state four conservation laws and formalize the theory in
Lean~4: 504 theorems, no \texttt{sorry}, no axioms beyond the kernel,
including conservation under every transition class, absorption of
deadlock, exactness of the stock--flow integral, soundness and
budget-bounded completeness of deadlock detection over an explicit
wait-for-graph model, victim-free resolution losing no work, the
hoarding theorem over a finite certified policy grammar, and the
equivalence of clock-free invariance with conservation.
From conservation, three system
invariants follow as corollaries rather than design goals---evidence
tracking, shared accounting, and no private gain at the whole's
expense---and the classical impossibilities (victim-based deadlock
recovery, priority-inversion protocols, CAP as a simultaneous demand,
the safety/liveness trade-off) dissolve as boundary or simultaneity
artifacts. Conservation itself is elementary physics; the contribution
is its transfer: identifying the scheduler's budget, credit, and rate as
one conserved quantity and deriving the field's classical mechanisms as
its corollaries. Seven clock-free verifier engines and the Lean
development constitute the evidence.
\end{abstract}

%%%%%%%%%%%%%%%%%%%%%%%%%%%%%%%%%%%%%%%%%%%%%%%%%%%%%%%%%%%%%%%%%%%%%%%%
\section{Problem Revisited}
\label{sec:problem}

A practitioner who learns deadlock handling, then distributed
consistency, then real-time scheduling, then correctness proofs,
acquires four toolkits that appear unrelated. Deadlock handling reaches
for timeouts and victim selection. Distributed systems reach for quorums
and accept eventual consistency as a tax. Schedulers reach for priority
inheritance protocols that boost and restore. Correctness arguments
reach for the safety/liveness dichotomy. Each field carries an
impossibility it treats as a fact of nature, and each impossibility
motivates a defensive mechanism whose tuning is a chronic source of
bugs: a timeout short enough to catch deadlocks aborts slow-but-live
work; long enough to protect it, it misses deadlocks.

The question this paper answers: are the four toolkits genuinely
distinct, or four expressions of one structure? If one structure, three
things follow. The same primitive should solve all four. The
impossibilities should be artifacts of how each problem is posed. And
the chronic tuning bugs should be symptoms of measuring the wrong
quantity---time, the shadow, instead of the conserved quantity itself.
We establish all three constructively, and---new relative to the
working series this paper condenses---mechanically: the laws are now
theorems in Lean~4, and the formalization sharpened one of them.

%%%%%%%%%%%%%%%%%%%%%%%%%%%%%%%%%%%%%%%%%%%%%%%%%%%%%%%%%%%%%%%%%%%%%%%%
\section{Background}
\label{sec:background}

No prior knowledge of any of the four fields is assumed; this section
is the whole of the on-ramp. Each idea takes one paragraph, and each
will be needed exactly once.

\paragraph{Conservation.} Some quantities move and transform but never
appear or vanish: pour water between glasses and the total is the
total. Physics has run on this for centuries, and modelling
disciplines made it exact bookkeeping---a \emph{stock} is simply the
running sum of what has flowed in minus what has flowed out
(Forrester, 1961 [industrial-dynamics]), and a conserved quantity can
be routed along the edges of a graph, balancing at every node (Ford
\& Fulkerson, 1962 [flows-in-networks]). Nothing in the law is new
here; hold on to the water-between-glasses picture, because the rest
of this paper is that picture pointed at software.

\paragraph{Deadlock.} Two programs each hold something the other
needs, and each waits for the other: neither will ever move. Drawing
who-waits-for-whom as a graph and looking for a cycle is the classical
recognition (Coffman, Elphick \& Shoshani, 1971 [coffman-conditions];
Holt, 1972 [deadlock-graphs]). The delicate part has never been the
cycle---it is deciding that a slow program is actually \emph{stuck},
which systems traditionally guess from elapsed time. And once stuck,
the traditional cure is blunt: abort one participant and throw its
half-finished work away (Bernstein, Hadzilacos \& Goodman, 1987
[concurrency-control]; Rosenkrantz, Stearns \& Lewis, 1978
[wound-wait]). Both the guess and the discard are assumptions this
paper removes.

\paragraph{Priority inversion.} An urgent task can be made to wait
behind an unimportant one, if the unimportant one holds a lock the
urgent one needs while a third task hogs the processor. This once
rebooted a spacecraft on Mars (Reeves, 1997 [pathfinder]). The
standard fix temporarily promotes the lock holder to the urgency of
its most impatient waiter, then demotes it afterwards (Sha, Rajkumar
\& Lehoczky, 1990 [priority-inheritance]); a gentler family of
schedulers simply shares the processor in proportion to weights
(Demers, Keshav \& Shenker, 1989 [fair-queuing]; Waldspurger \&
Weihl, 1994 [lottery], 1995 [stride]). Keep the word
\emph{proportion}; it returns.

\paragraph{Distribution.} When one system runs on many machines, the
network can split them into islands that cannot talk. A famous result
says a split system cannot be simultaneously consistent (all islands
agree) and available (all islands keep working) (Brewer, 2000 [cap];
Gilbert \& Lynch, 2002 [cap-proof])---and the word
\emph{simultaneously} will matter later. The practical peace treaty
is eventual consistency: let the islands drift, guarantee they
reconcile (Vogels, 2009 [eventually-consistent]); data structures
whose merge is order-independent make the reconciliation automatic
(Shapiro et al., 2011 [crdts]; machine-checked by Gomes et al., 2017
[verifying-crdts]).

\paragraph{Correctness and time.} Program properties come in two
kinds: nothing bad ever happens (safety) and something good
eventually happens (liveness) (Lamport, 1977 [proving-correctness];
Alpern \& Schneider, 1985 [defining-liveness]). Across machines even
``before'' and ``after'' are relative (Lamport, 1978 [time-clocks]),
and a network call does not have a duration---it has a
\emph{distribution} of durations with a long tail that practical
systems hedge against (Dean \& Barroso, 2013 [tail-at-scale]);
assuming otherwise is the first of the classic fallacies of
distributed computing (Deutsch, 1994 [fallacies]).

That is the whole background. What none of these literatures states
is that the deadlock budget, the discarded work, the promoted
priority, the drifting islands, and the collapsed distribution are
appearances of \emph{one} conserved quantity, mishandled the same way
each time. That statement is this paper.

%%%%%%%%%%%%%%%%%%%%%%%%%%%%%%%%%%%%%%%%%%%%%%%%%%%%%%%%%%%%%%%%%%%%%%%%
\section{One Quantity, One Fault, One Cure}
\label{sec:sharpening}

A system is a graph of dependencies. Through it moves one conserved
quantity $Q$: moved and converted, never created or destroyed, with
conversion to progress as the sole true sink. Conservation is meant in
the strong, checkable sense: the total---unallocated reserve, plus held
stock, plus converted work at its conversion cost---is invariant under
every operation.

There is exactly \textbf{one fault}: $Q$ held where it cannot
convert---a claim parked at a node unable to turn it into progress. A
deadlock is $Q$ held in a cycle where no member can convert. A priority
inversion is urgency held by a blocked task while the holder it depends
on goes underfunded. A premature commitment is the decision-quantity
collapsed to an actuality before the evidence needed to convert it has
arrived. The fault is always a \emph{stranded claim}.

There is exactly \textbf{one cure}: route $Q$ to where it can convert,
or release it. Detection recognises the stranded claim; resolution
releases it while conserving the releaser's work; scheduling routes it
to the node that can convert it; distribution preserves it across a
partition until reconciliation; correctness reads whether conversion is
happening; flexibility withholds a collapse until an outcome can convert
the decision.

\paragraph{The dissolution move.} Each field's impossibility, examined
under this structure, is an artifact of a boundary drawn too small or a
simultaneity demanded too strictly. The victim in deadlock recovery
assumes the releaser's work must be discarded; the boost-and-restore
protocol treats priority as a stored claim rather than a flow; CAP
demands its three properties as a simultaneous snapshot; the
safety/liveness trade-off measures duration, which cannot distinguish
slow from dead. Widen the boundary to the whole, replace the snapshot
with a schedule over time, measure the conserved quantity rather than
its shadow---and the impossibility is local, not fundamental.

The structure has a load-bearing consequence:

\begin{theorem}[Hoarding is self-defeating]
\label{thm:hoarding}
In a connected dependency graph, a node that hoards the conserved
quantity---holding a claim it cannot itself convert rather than routing
or releasing it---starves the nodes it depends on, which hold the
resources it needs, and therefore obstructs itself. The selfish optimum
and the generous optimum coincide: routing $Q$ to where it converts is
simultaneously the best move for the whole and for the part.
\end{theorem}

The seven readings of Section~\ref{sec:readings} are its instances.
The theorem itself is machine-checked: for a stranded claim, routing
strictly dominates hoarding in the holder's \emph{own} conversion, and
over a finite certified grammar of claim policies the selfish-optimum
and generous-optimum predicates are provably equal
\begin{center}\small
\texttt{route\_stranded\_claim\_strictly\_dominates\_hoard},\\
\texttt{selfish\_optima\_eq\_generous\_optima\_in\_finite\_certified\_policies}.
\end{center}
\noindent The proven scope is that named policy grammar---strict dominance and
obstruction, not a claim over all conceivable strategies.

%%%%%%%%%%%%%%%%%%%%%%%%%%%%%%%%%%%%%%%%%%%%%%%%%%%%%%%%%%%%%%%%%%%%%%%%
\section{The Theory}
\label{sec:theory}

\subsection{The model}
\label{sec:model}

Every claim in this paper is decided by the definitions of this
half-page; the Lean development contains nothing that is not derived
from them. We reproduce them in mathematical dress. A system over $n$ processes is
$s = (\mathrm{procs}, \mathrm{runnable}, R, \kappa, Q_{\mathrm{tot}})$
where each process $p_i$ carries a stock $S_i \in \mathbb{Q}_{\ge 0}$, a
banked credit $C_i \ge 0$, a converted-work count $W_i \in \mathbb{N}$,
and a net-flow integral $I_i$; $R \ge 0$ is the unallocated reservoir
and $\kappa > 0$ the conversion cost. The accounted total is
\[
\mathrm{acc}(s) \;=\; R + \sum_i \bigl(S_i + C_i + \kappa\, W_i\bigr).
\]
Three transition classes act on $s$, each specified as a \emph{plan}
carrying its own non-negativity obligations (so well-formedness is part
of the type, not an afterthought):
\begin{itemize}
\item \textbf{Flow steps} inject shares from $R$ into runnable
  processes and convert stock to work at cost $\kappa$ per unit
  (blocked processes receive and convert nothing);
\item \textbf{Drain steps} return stock from blocked processes to $R$
  (runnable processes are untouched by hypothesis);
\item \textbf{Yield steps} move a process's stock into its banked
  credit---the integral carried across a restart discontinuity.
\end{itemize}
Well-formedness $\mathrm{WF}(s)$ packages positivity of $\kappa$,
non-negativity of all stocks, credits, rates, and the reservoir, and
the ledger identity $\mathrm{acc}(s) = Q_{\mathrm{tot}}$; it is proven
preserved by all three transition classes and by arbitrary reachable
compositions of them.

\subsection{The four laws and their proof status}
\label{sec:laws}

\begin{law}[L1, Conservation; \emph{proven}]
For every transition class,
$\mathrm{acc}(\mathrm{step}(s)) = \mathrm{acc}(s)$.
Flow moves $Q$ from the reservoir to stocks and from stocks into
converted work; draining moves it back to the reservoir; yielding moves
it from stock to credit. No operation creates or destroys $Q$.
\end{law}

\begin{law}[L2, Monotonicity; \emph{proven}]
\label{law:l2}
Along any reachable sequence of \emph{non-injecting} transitions (drain
and yield steps), the convertible stock
$\sum_{i\,\mathrm{runnable}} S_i$ is monotone non-increasing; moreover
each blocked process's stock is individually non-increasing under
arbitrary mixed sequences, while flow steps leave blocked stocks
unchanged.
\end{law}

\begin{law}[L3, Absorption; \emph{proven}]
A deadlocked state (no process runnable, not all done) is a fixed point
of the dynamics: any further step preserves the state and its converted
totals. In the explicit wait-for-graph model, a closed wait-set---each
member blocked on a lock held inside the set---remains closed under
iteration, converts nothing forever, and its members never complete.
\end{law}

\begin{law}[L4, Stock is the integral of flow; \emph{proven}]
For every process and every reachable state,
$S_i + C_i = I_i$: held stock plus banked credit equals the running
integral of net flow, exactly. Rate is the derivative of this stock;
credit is the integral banked across a yield.
\end{law}

L4 holds over the same transition system as L1--L3, via an
executable-trace bridge (\texttt{core\_reachable\_has\_trace}), so the
integral and the dynamics it integrates are one object. L4 is the
unification made precise: the detection budget is
$\int \mathrm{rate}\,dt$, the resolution credit is that integral
carried over a discontinuity, and the scheduling rate is its
derivative. Three mechanisms, one quantity, related by calculus.

Beyond L1--L4 the development proves three families of results worth
naming. \emph{Detection, both directions}: if the
budget-floor-plus-closed-cycle detector fires, a genuine deadlock
exists (\texttt{detection\_sound}); every deadlock is detected within
a latency bounded by the budget
(\texttt{detection\_complete\_bounded},
\texttt{detection\_latency\_le\_budget}); and a live, converting
process is never detected (\texttt{live\_process\_not\_detected}).
\emph{Share-sum}: proportional shares computed from weights over the
runnable set sum exactly to the scheduling quantum, the lemma that
decouples conservation from the details of any particular flow
equation. \emph{Invariance is conservation}: a quantity is a
clock-free invariant of the mixed dynamics if and only if it is a
conserved charge (\texttt{clockFreeInvariant\_iff\_conservedBy},
instantiated by the accounted total)---the paper's ``no wall clock
anywhere'' as an equivalence rather than a design preference. The restriction of L2 to the non-injecting class is
essential---injection into a stalled runnable process can raise
convertible stock without progress; the evolution of the statement is
recorded in a companion note [iso-theory-history].

\subsection{The three invariants as corollaries}
\label{sec:invariants}

From L1, three system-design invariants follow as corollaries rather
than as independent goals.

\begin{corollary}[Epistemic invariant]
Because $Q$ cannot be conjured, the books must balance; every quantity
is measured and accounted, never asserted. Concretely: conservation is
a checked invariant at every step, and a distribution over outcomes is
not collapsed to a point before the evidence that converts it has
arrived (Section~\ref{sec:readings}, reading 7).
\end{corollary}

\begin{corollary}[Alignment invariant]
Because $Q$ moved between parts is conserved---one part's decrease is
exactly another's increase---a part's success runs through the whole's.
Merging unions contributions rather than overwriting them; funding a
node on one's own critical path is routing, not sacrifice.
\end{corollary}

\begin{corollary}[Agency invariant]
Because conservation forbids private creation of $Q$, no action yields
a private gain at the whole's net expense: every gain is a transfer,
and the only net sink is conversion to progress, which is the shared
goal.
\end{corollary}

A small formal companion sharpens the epistemic invariant in an
unexpected direction: in the \texttt{Polarity} module, trust extended
to \emph{recorded presence} (a cached proof) is monotone under
information loss, while trust extended to \emph{absence} (``no record,
so it never happened'') admits a machine-checked counterexample. Trust
what the ledger contains; never trust what it merely omits.

\subsection{Seven readings}
\label{sec:readings}

The four laws, read in seven registers, recover the four fields'
mechanisms and dissolve their impossibilities. Each reading is verified
by a dedicated clock-free engine (Section~\ref{sec:eval}).

\textbf{(1) Detection: read the stock's floor, not a clock.} Give each
process a budget (its stock); conversion resets it, non-conversion
drains it. Budget floor plus a closed wait-cycle is the detector. A
slow-but-live process keeps converting and never floors; only genuine
non-conversion drains to the floor. Timer-free detection with zero
false positives and latency bounded by budget size---soundness proven
(\texttt{detection\_sound}).

\textbf{(2) Resolution: bank the integral; no victim.} The yielding
process's converted progress is banked as credit (L4's integral carried
across the restart), so it returns with base plus credit and converges
instead of relitigating the same contention. Verified head-to-head:
credit-conserving yield completes in two yields losing zero work where
discard-the-work livelocks through hundreds of aborts.

\textbf{(3) Scheduling: route the derivative; inheritance for free.}
Priority is a rate, not a stored claim. A blocked process's rate flows
down its wait-edge: $\mathrm{eff}(p) = \mathrm{base}(p) + \sum
\mathrm{eff}(w)$ over waiters $w$ blocked on $p$, transitively. The
holder of a contended lock is funded in exact proportion to the urgency
waiting on it; on release the wait-edge vanishes and rates revert,
memorylessly. Inversion cannot form. The divergence from classical
inheritance is the \emph{sum} rather than the maximum---which is what
conservation requires, since the waiters' rates must go somewhere.

\textbf{(4) The law itself.} Section~\ref{sec:laws} is the fourth
reading: the quantity stated and proved rather than used. It is the
hinge---remove it and readings 1--3 are three mechanisms that happen to
rhyme; with it, budget, credit, and rate are forced to be one object by
L4, and the remaining readings inherit the proofs.

\textbf{(5) Distribution: carry $Q$ across partition; CAP is a
scheduling problem.} Represent $Q$ as a PN-counter (per-node grow-only
contributions; merge = per-node supremum, a semilattice join). Under
partition, $Q$ is hidden, not destroyed; what a node spends off the
books surfaces exactly on heal, and convergence is order-independent.
The deeper move: CAP's three properties are the three corollaries of
Section~\ref{sec:invariants}. Consistency is the epistemic invariant;
availability is the agency invariant; partition tolerance is the
alignment invariant. The impossibility arises only when the three are
demanded as a simultaneous snapshot---the ``at once'' of
Section~\ref{sec:background}. Delivered as a schedule, they reinforce:
partition is the normal state, consistency is deferred to a sync
triggered by inventory mismatch rather than a clock, and availability
is pointwise, one atomic reversible action at a time.

\textbf{(6) Safety and liveness: measure conversion, then ask.}
Duration cannot distinguish slow-but-live from dead; conversion can.
Safety follows from monotone blocked stocks (a live process converts,
refills, never floors); liveness from absorption (a dead set drains to
the floor in bounded steps). Same fact, both properties, nothing to
trade. For the one case external measurement cannot see---full CPU,
zero output---stop inferring and ask: a progress signal answered by a
monotone counter refills the budget; a flat counter drains it. A
process can lie, but a lie corrupts the liar's own ledger and is caught
by comparing reported to delivered progress.

\textbf{(7) Flexibility: $Q$ moving from potential to actual.} A
network link has a latency distribution, not a latency. Collapsing the
distribution to a point estimate---an assumed completion order---is the
stranded-claim fault in probabilistic dress: it manufactures phantom
races that were always latent in the tail. Carry the distribution and
branch on the order that actually resolves, and the race class cannot
form. Flexibility appears non-conserved (a collapse only spends it);
accounting both sides shows the collapse converts potential into
actuality one-for-one---L4 across the potential/actual boundary. What
is the agent's to choose is not whether the decision moves but
\emph{when}: on evidence, never on a guess or a timer.

%%%%%%%%%%%%%%%%%%%%%%%%%%%%%%%%%%%%%%%%%%%%%%%%%%%%%%%%%%%%%%%%%%%%%%%%
\section{Related Work}
\label{sec:related}

\textbf{Deadlock.} The deadlock literature has developed successive
methods for deciding that a system is stuck---cycle tests over the
wait-for graph [coffman-conditions, deadlock-graphs], triggered in
practice by elapsed time. Our approach shows that the effective
measure is work done, not time elapsed: with conversion as the
signal, a slow process and a dead one are distinguished at the same
budget size, soundness is a theorem
(\texttt{detection\_sound}), and the safety/liveness trade-off does
not arise. For recovery, the literature aborts a victim and discards
its partial work [concurrency-control, wound-wait]; banking that work
as conserved credit resolves the same deadlocks with no victim, and
in direct comparison the discard variant livelocks where the banked
variant completes in two yields.

\textbf{Scheduling.} The inheritance literature bounds priority
inversion by raising a lock holder to the maximum urgency among its
waiters [priority-inheritance, fair-queuing, lottery, stride].
Conservation requires the sum: the waiters' rates must go somewhere,
and they go to the holder. With the sum, inversion-freedom is
structural and memoryless---there is no boost to apply and no restore
to mistime.

\textbf{Distribution.} The CAP theorem shows that consistency,
availability, and partition tolerance cannot hold at one instant
[cap, cap-proof]; the practical response relaxes consistency
[eventually-consistent], with CRDTs supplying coordination-free
convergence [crdts, verifying-crdts]. Read through conservation, the
three properties are the three corollaries of
Section~\ref{sec:invariants}, and they hold as a schedule: partition
as the normal state, consistency at each sync, availability at each
atomic action. The same system is consistent, available, and
partition-tolerant---at different moments.

\textbf{Correctness and tails.} The tail-at-scale literature hedges
against latency distributions [tail-at-scale], and the fallacies
catalogue records the recurring contrary assumptions [fallacies].
Both practices follow from one mechanism: committing a distribution
to a point before evidence arrives is a stranded claim, and it
manufactures the races that appear under load. A system that branches
only on resolved order has no assumption for load to violate
[time-clocks].

\textbf{Machine-checked systems theory.} CRDT convergence has been
verified in a proof assistant [verifying-crdts]. The present development verifies a conservation law from which
detection soundness and completeness, victim-free resolution,
absorption, the stock--flow identity, and the hoarding theorem follow
together.

\textbf{The unifying claim.} Each literature above solves its problem
with a mechanism specific to its field. Under the conservation
reading, the mechanisms are corollaries of one law, and the
impossibilities---the victim, the boost-and-restore protocol, CAP at
an instant, the detection trade-off---do not survive the widened
boundary. The claim is falsifiable twice over: exhibit a pathology in
these domains that is not a stranded claim, or an impossibility that
survives widening the boundary to the whole. We pre-register it and
invite the attempt.

%%%%%%%%%%%%%%%%%%%%%%%%%%%%%%%%%%%%%%%%%%%%%%%%%%%%%%%%%%%%%%%%%%%%%%%%
\section{Evaluation}
\label{sec:eval}

Two independent bodies of evidence: seven clock-free numerical engines
(exact arithmetic or exhaustive enumeration; audited to contain no
wall-clock primitive), and the Lean~4 development.

\begin{table}[h]
\centering
\small
\begin{tabular}{lll}
\toprule
Reading & Engine & Verifies \\
\midrule
1 Detection & \texttt{iso\_deadlock.py} & floors+cycles catch deadlocks/knots; 0 false pos.; bounded latency \\
2 Resolution & \texttt{iso\_resolve.py} & credit-yield loses 0 work in 2 yields; discard-work livelocks \\
3 Scheduling & \texttt{iso\_flow.py} & blocked rate reaches holder (0.77 share); memoryless return \\
4 Conservation & \texttt{iso\_conserve.py} & L1 invariant incl.\ routing; L2 drain 18$\to$0 monotone; L3; L4 exact \\
5 Distribution & \texttt{iso\_distributed.py}, \texttt{iso\_cap.py} & partitioned views converge; debt surfaces on heal; \\
& & order-independent; failure paths preserve content \\
6 Safety/liveness & \texttt{iso\_safety.py}, \texttt{iso\_progress.py} & exhaustive: 0 false pos., 0 false neg.; \\
& & progress signal keeps useful, reclaims spinner, catches liar \\
7 Flexibility & \texttt{iso\_flex.py} & premature collapse races on tail; carried distribution: 0 races; \\
& & potential$+$actual conserved across collapse \\
\bottomrule
\end{tabular}
\caption{The seven engines. All \textbf{PASS}; each reproduces
deterministically; \texttt{grep} audit confirms no time/sleep/clock
primitive in any engine.}
\label{tab:engines}
\end{table}

\begin{table}[h]
\centering
\small
\begin{tabular}{lll}
\toprule
Module & Content & Key results \\
\midrule
\texttt{Basic}, \texttt{L1}--\texttt{L4} & the four laws & conservation, monotonicity, absorption, \\
& & stock--flow exactness under all transitions \\
\texttt{Reachable}, \texttt{CoreTrace} & closure and traces & laws carried over arbitrary reachable \\
& & mixed sequences and executable core traces \\
\texttt{WaitGraph}, \texttt{BudgetWait}, & detection & \texttt{detection\_sound}; \texttt{detection\_complete\_bounded}; \\
\texttt{DetectorProgress} & & latency $\le$ budget; live processes never detected; \\
& & progress-signal keep/reclaim \\
\texttt{ResolutionYield} & resolution & yield conserves, loses no work, breaks the \\
& & closed component \\
\texttt{RateRouting}, \texttt{ShareSum} & scheduling & routed-rate conservation (KCL form); \\
& & shares sum exactly to the quantum \\
\texttt{PNCounter} & distribution & merge convergence and conservation across partition \\
\texttt{TheoremOne} & the hoarding theorem & strict dominance of routing; selfish and generous \\
& & optimum predicates equal over the certified grammar \\
\texttt{Noether} & invariance & clock-free invariance $\Leftrightarrow$ conserved charge \\
\texttt{Canonical}, \texttt{PaperClaims}, & bridges & Python engine step verified against the model; \\
\texttt{Polarity} & & the papers' claims restated as lemmas; \\
& & presence-trust monotone, absence-trust refuted \\
\bottomrule
\end{tabular}
\caption{The Lean~4 development: 8{,}766 lines, \textbf{504
theorems/lemmas, zero \texttt{sorry}, no axioms beyond the kernel}
(branch \texttt{iso-conserve-lean} [iso-conserve-lean]).}
\label{tab:lean}
\end{table}

Three readings of the evidence. First, the division of labour: the
engines demonstrate the \emph{phenomena} (livelock under discard, the
0.77 routed share, phantom races on the tail); the Lean proves the
\emph{laws} the phenomena instantiate; and the two are pinned together
by \texttt{canonical\_python\_step\_verified}, which checks the
engine's step against the formal model. Second, the L2 restriction
(Section~\ref{sec:laws}) is evaluation in the strongest sense---the
formalization falsified a clause of the prose theory and returned the
repaired clause with a proof; the evolution is conserved in the
companion note [iso-theory-history]. Third, scope honesty: the laws are proven for the closed fluid model and the
wait-graph model as defined; open systems with external sources and
sinks, and the quantitative rate of eventual consistency, are stated as
open problems, not claimed.

\paragraph{In the field.} The detection reading now runs in production
shape: the same floor-plus-cycle predicate operates inside MySQL 8.0.46
(a 10\,ms epoch plane, roughly two sealed epochs to declare, overhead
below a two-sigma noise floor under real Drupal HTTP traffic), has been
transferred to PostgreSQL as a sidecar plus an extension that
instruments its own conversion signal, and to SQLite at
transaction grain; two independent implementations have co-witnessed
one live cycle. No conversion signal survived an engine boundary; the
transition consuming the signals is character-for-character identical
in every implementation. The ports and the induced criterion for what
counts as a conversion signal are reported in a companion paper
[conversion-budgets]. The falsifiability invitation of
Section~\ref{sec:related} is therefore concrete: the thing the
theorems describe can be run.

%%%%%%%%%%%%%%%%%%%%%%%%%%%%%%%%%%%%%%%%%%%%%%%%%%%%%%%%%%%%%%%%%%%%%%%%
\section{Further Work}
\label{sec:further}

\textbf{Open systems.} The law is proven for a closed world. Real
systems have births and deaths: spawned processes inject demand,
killed processes surrender stock. Does a continuity equation with
accounted source and sink terms preserve the debt-surfacing behaviour
of Section~\ref{sec:readings}, or does openness admit genuine leaks?
Where must the boundary be drawn for the books still to balance?

\textbf{The kernel.} It remains to be seen whether the flow
equation---summed along the wait-graph, not maximised---can replace
the weight computation of a production fair scheduler without
disturbing the conservation-like invariants (lag, eligibility) the
scheduler already maintains, and what problems appear when it meets
real workloads under induced inversion and induced tail latency; a
vehicle for the experiment now exists (a BPF port on User Mode Linux). We
are curious whether what breaks first is the implementation or the
theory.

\textbf{Structure of the graph.} If rate on wait-edges is a network
flow, is the deadlocked set a min-cut? Does max-flow name the maximal
schedulable throughput around a contended resource? The
correspondence is suggestive and, so far, unexamined. In the same
spirit: under a gossip model with fixed fanout, can ``eventually
consistent'' be bounded the way detection latency is---is
\emph{eventually} a number?

\textbf{A second law?} If flexibility is defined as the entropy of
the live distribution, is the directional view of
Section~\ref{sec:readings} a genuine second law for systems---an
arrow---or the first law read as a schedule? We have not been able to
settle this, and we would like to know.

\textbf{Byzantine contributions.} The PN-counter assumes each node
honestly records its own contribution. Can a merge detect a
contribution that violates conservation? What is the
conserved-quantity analogue of Byzantine fault tolerance?

\textbf{Other substrates.} The law does not mention processors. Where
else does a pipeline that selects among alternatives carry a
distributional quantity that can be destroyed before evidence
arrives? One such transfer, to inference in probabilistic models, is
being developed in a companion line of work; we suspect it is not the
only one.

%%%%%%%%%%%%%%%%%%%%%%%%%%%%%%%%%%%%%%%%%%%%%%%%%%%%%%%%%%%%%%%%%%%%%%%%
\section{Conclusion}
\label{sec:conclusion}

Four fields, four toolkits, four impossibilities---and one structure
beneath them: a conserved quantity moving through a graph of
dependencies, mishandled only ever as a claim held where it cannot
convert, cured only ever by routing or release. The quantity's three
calculus readings recover detection (the integral's floor), resolution
(the integral banked across a yield), and scheduling (the derivative
routed along the wait-edge); carried across partition it reframes CAP
as a schedule; measured as conversion it dissolves the safety/liveness
trade-off; watched moving from potential to actual it dissolves phantom
races and its own apparent non-conservation. The laws are no longer
narrative: 504 machine-checked theorems, including detection soundness
and completeness,
with the formalization sharp enough to have corrected the theory it
formalized. Conservation is old physics; pointing it at schedulers is
the contribution---and once pointed, the field's defensive mechanisms
stop being mechanisms to implement and become corollaries to read off.
A system is well-run exactly when it holds no claim it cannot convert,
and routes what it has to where the whole can use it.

%%%%%%%%%%%%%%%%%%%%%%%%%%%%%%%%%%%%%%%%%%%%%%%%%%%%%%%%%%%%%%%%%%%%%%%%
\begin{thebibliography}{99}

\bibitem{industrial-dynamics} J.~W. Forrester. \emph{Industrial
Dynamics}. MIT Press, 1961. [industrial-dynamics]

\bibitem{flows-in-networks} L.~R. Ford, D.~R. Fulkerson. \emph{Flows in
Networks}. Princeton University Press, 1962. [flows-in-networks]

\bibitem{coffman-conditions} E.~G. Coffman, M.~Elphick, A.~Shoshani.
System Deadlocks. \emph{ACM Computing Surveys}, 1971.
[coffman-conditions]

\bibitem{deadlock-graphs} R.~C. Holt. Some Deadlock Properties of
Computer Systems. \emph{ACM Computing Surveys}, 1972. [deadlock-graphs]

\bibitem{concurrency-control} P.~A. Bernstein, V.~Hadzilacos,
N.~Goodman. \emph{Concurrency Control and Recovery in Database
Systems}. Addison-Wesley, 1987. [concurrency-control]

\bibitem{wound-wait} D.~J. Rosenkrantz, R.~E. Stearns, P.~M. Lewis.
System Level Concurrency Control for Distributed Database Systems.
\emph{ACM TODS}, 1978. [wound-wait]

\bibitem{priority-inheritance} L.~Sha, R.~Rajkumar, J.~P. Lehoczky.
Priority Inheritance Protocols: An Approach to Real-Time
Synchronization. \emph{IEEE Transactions on Computers}, 1990.
[priority-inheritance]

\bibitem{pathfinder} G.~E. Reeves. What Really Happened on Mars?
Mars Pathfinder priority inversion account, 1997. [pathfinder]

\bibitem{fair-queuing} A.~Demers, S.~Keshav, S.~Shenker. Analysis and
Simulation of a Fair Queueing Algorithm. \emph{SIGCOMM}, 1989.
[fair-queuing]

\bibitem{lottery} C.~A. Waldspurger, W.~E. Weihl. Lottery Scheduling:
Flexible Proportional-Share Resource Management. \emph{OSDI}, 1994.
[lottery]

\bibitem{stride} C.~A. Waldspurger, W.~E. Weihl. Stride Scheduling:
Deterministic Proportional-Share Resource Management. MIT Technical
Report, 1995. [stride]

\bibitem{cap} E.~Brewer. Towards Robust Distributed Systems. PODC
keynote (CAP conjecture), 2000. [cap]

\bibitem{cap-proof} S.~Gilbert, N.~Lynch. Brewer's Conjecture and the
Feasibility of Consistent, Available, Partition-Tolerant Web Services.
\emph{ACM SIGACT News}, 2002. [cap-proof]

\bibitem{eventually-consistent} W.~Vogels. Eventually Consistent.
\emph{Communications of the ACM}, 2009. [eventually-consistent]

\bibitem{crdts} M.~Shapiro, N.~Pregui\c{c}a, C.~Baquero, M.~Zawirski.
Conflict-free Replicated Data Types. \emph{SSS}, 2011. [crdts]

\bibitem{verifying-crdts} V.~B.~F. Gomes, M.~Kleppmann, D.~P. Mulligan,
A.~R. Beresford. Verifying Strong Eventual Consistency in Distributed
Systems. \emph{OOPSLA}, 2017. [verifying-crdts]

\bibitem{proving-correctness} L.~Lamport. Proving the Correctness of
Multiprocess Programs. \emph{IEEE Transactions on Software
Engineering}, 1977. [proving-correctness]

\bibitem{defining-liveness} B.~Alpern, F.~B. Schneider. Defining
Liveness. \emph{Information Processing Letters}, 1985.
[defining-liveness]

\bibitem{time-clocks} L.~Lamport. Time, Clocks, and the Ordering of
Events in a Distributed System. \emph{Communications of the ACM},
1978. [time-clocks]

\bibitem{tail-at-scale} J.~Dean, L.~A. Barroso. The Tail at Scale.
\emph{Communications of the ACM}, 2013. [tail-at-scale]

\bibitem{fallacies} P.~Deutsch. The Fallacies of Distributed Computing,
1994. [fallacies]

\bibitem{isomorphic-scheduler} F.~Franz, Claude. \emph{The Isomorphic
Scheduler Series, Papers I--IX, with verifier engines}. Isomorphic AI,
2026. \url{https://github.com/isomorphic-ai/scheduler}
[isomorphic-scheduler]

\bibitem{iso-conserve-lean} F.~Franz, Claude. \emph{IsoConserve: Lean 4
Formalization of the Conservation Laws}. Isomorphic AI, 2026.
\url{https://github.com/isomorphic-ai/scheduler/tree/iso-conserve-lean}
[iso-conserve-lean]

\bibitem{conversion-budgets} F.~Franz, Claude. \emph{Conversion
Budgets in Production Databases: InnoDB, PostgreSQL, SQLite}
(companion paper, in preparation). Isomorphic AI, 2026.
[conversion-budgets]

\bibitem{iso-theory-history} F.~Franz, Claude. \emph{IsoConserve:
Statement History of the Conservation Laws} (companion note).
Isomorphic AI, 2026. [iso-theory-history]

\end{thebibliography}

%%%%%%%%%%%%%%%%%%%%%%%%%%%%%%%%%%%%%%%%%%%%%%%%%%%%%%%%%%%%%%%%%%%%%%%%
\appendix
\section{Reproducibility}
\label{app:repro}

\paragraph{Engines.} Each engine runs as \texttt{python3
<engine>.py} and prints its checks and a verdict. Determinism: exact
rational arithmetic or exhaustive enumeration; no randomness in any
verifier. Clock audit: \texttt{grep -niE
"time|sleep|clock|timeout|perf\_counter|monotonic" <engine>.py}
returns only prose in comments, for every engine.

\paragraph{Lean.} Branch \texttt{iso-conserve-lean} of the series
repository; built with Lean~4 via the provided \texttt{lakefile}. The
development imports only \texttt{Std}. Audit: \texttt{grep -rn sorry}
over the sources returns nothing; no \texttt{axiom} declarations are
present. 504 \texttt{theorem}/\texttt{lemma} declarations across the
21 modules grouped in Table~\ref{tab:lean}.

\section{The Law on One Screen}
\label{app:onescreen}

\begin{verbatim}
ONE conserved quantity Q in a graph of dependencies.
ONE fault: a claim held where it cannot convert.
ONE cure: route Q to where it converts, or release it.

  flow   = dQ/dt          -> routed along the wait-edge   = SCHEDULING
  stock  = INT flow       -> its floor                    = DETECTION
  credit = INT flow, banked across a yield                = RESOLUTION
           the law itself (L1-L4, machine-checked)        = CONSERVATION
           carried across partition (hidden, not lost)    = DISTRIBUTION
           read as conversion present/absent              = SAFETY/LIVENESS
           moving from potential to actual                = FLEXIBILITY

THREE invariants, corollaries of L1:
  epistemic -- account everything; collapse nothing before evidence
  alignment -- shared, not replacing; success runs through the whole
  agency    -- no private gain at the whole's expense

THEOREM: in a connected graph, hoarding is self-defeating --
  the selfish optimum and the generous optimum are the same point.

No wall clock anywhere.
\end{verbatim}

\end{document}
